\chapter[Group (Co)homology and Lie Cohomology]
  {Group (Co)homology and\\Lie Algebra Cohomology}
\label{ch:group-and-lie}

Throughout this chapter \(k\) is a commutative ring. Group modules are left
modules over the group algebra, and all groups under consideration are
discrete. For a Lie algebra \(\mathfrak g\) over \(k\), a coefficient
module is a left module over \(U(\mathfrak g)\), or equivalently a
\(k\)-module with a compatible \(\mathfrak g\)-action. The trivial module
\(k\) is defined by the augmentation of the respective algebra. We
apply the resolution arguments of Chapter~\ref{ch:derived-functors} and
the composition spectral sequence to these two augmented algebras.
The final section constructs coefficient group homology from a right
bar resolution and identifies it with derived coinvariants.

\section{Invariants and the bar resolution}

\begin{proposition}[Invariants and group cohomology]\label{res:A01}
For a group \(G\) and a left \(k[G]\)-module \(M\), evaluation at \(1\)
identifies \(\operatorname{Hom}_{k[G]}(k,M)\) with \(M^G\), naturally in
\(M\). Its right derived functors, written \(H^n(G,M)\), are
\(\operatorname{Ext}_{k[G]}^n(k,M)\), computed either from an injective
resolution of \(M\) or a projective resolution of \(k\).
\end{proposition}

\begin{proof}
The image of \(1\) under a module map is fixed by each \(g\), and a
fixed element \(m\) defines the map \(a\mapsto am\). This is a natural
identification of left exact functors to \(k\)-modules. The category
of modules over \(k[G]\) has enough injectives and projectives.
Theorems~\ref{res:R09} and~\ref{res:R11} give the derived
identifications and their independence of resolution.
See \cite[Section 6.1]{weibel}.
\end{proof}

\begin{theorem}[The bar resolution and nonhomogeneous cochains]
\label{res:A02}
Let \(B_n(G)=k[G^{n+1}]\), the free \(k\)-module on \((n+1)\)-tuples,
with the diagonal left \(G\)-action. With the alternating deletion
boundary and the augmentation \(B_0(G)\to k\) sending every vertex
to \(1\), this is a free \(k[G]\)-resolution of \(k\). Put
\(C^n(G,M)=\operatorname{Map}(G^n,M)\), with \(C^0(G,M)=M\).
The Hom complex from this resolution identifies with these cochains.
Its differential in degree \(n\) is
\begin{align*}
 (df)(g_1,\ldots,g_{n+1})
  &=g_1 f(g_2,\ldots,g_{n+1})\\
  &\quad+\sum_{i=1}^{n}(-1)^i
          f(g_1,\ldots,g_i g_{i+1},\ldots,g_{n+1})\\
  &\quad+(-1)^{n+1} f(g_1,\ldots,g_n).
\end{align*}
In particular \(H^n(G,M)\) is its \(n\)-th cohomology.
\end{theorem}

\begin{proof}
Two deleted positions in the boundary of a boundary occur in two
orders with opposite signs. The degree-one boundary has equal
augmentation on its two terms, so the augmented differential squares to
zero everywhere. Each tuple has a unique representative with its first
entry \(1\) after translation, which makes \(B_n(G)\) free as a left
\(k[G]\)-module. On the underlying augmented \(k\)-complex set
\[
 s_{-1}(1)=(1),\qquad
 s_n(g_0,\ldots,g_n)=(1,g_0,\ldots,g_n).
\]
Deleting the first entry of \(s_n\) gives the original tuple.
Deleting any other entry occurs with sign opposite to the corresponding
term of \(s_{n-1}\partial\). This proves \(\partial s+s\partial=1\),
including degrees zero and \(-1\). The homotopy need only be
\(k\)-linear: the resulting exactness of the underlying sequence is
exactness of the sequence of \(k[G]\)-modules.

The notation \([g_1|\cdots|g_n]\) denotes the tuple
\((1,g_1,g_1g_2,\ldots,g_1\cdots g_n)\) in a basis over \(k[G]\).
Deleting the first vertex in such a tuple gives
\(g_1[g_2|\cdots|g_n]\). Deleting the \(i\)-th internal vertex
multiplies its two adjacent increments, and deleting the last
vertex drops \(g_n\). Thus precomposition by the boundary gives
the displayed cochain formula. This uses the unsigned
precomposition convention for the Hom from a chain resolution,
as in Theorem~\ref{res:R11}. Equivalently, a homogeneous
cochain is a map \(\varphi:G^{n+1}\to M\) satisfying
\(\varphi(a\boldsymbol g)=a\varphi(\boldsymbol g)\).
Its values on the displayed basis determine it uniquely. The
identification commutes with all coefficient maps, and the free
resolution computes cohomology by Proposition~\ref{res:A01}.
See \cite[Sections 6.5.1--6.5.6]{weibel}.
\end{proof}

\begin{corollary}[Low degrees and the coefficient sequence]
\label{res:A03}
There are natural identifications
\[
 H^0(G,M)=M^G,\qquad
 H^1(G,M)=
 \frac{\{f:G\to M\mid f(gh)=f(g)+g f(h)\}}
      {\{g\mapsto g m-m\mid m\in M\}} .
\]
A short exact sequence of coefficients yields a natural long exact
sequence with connecting map given by the differential of a lifted cochain.
\end{corollary}

\begin{proof}
At degrees zero and one the differential in
Theorem~\ref{res:A02} gives precisely the indicated conditions.
The functors \(\operatorname{Hom}_{k[G]}(B_n(G),-)\) are exact because
the \(B_n(G)\) are free. The coefficient sequence thus gives a short
exact sequence of cochain complexes. If a cocycle in the quotient
complex is lifted to the middle complex, its differential belongs to
the kernel complex. A second lift changes this differential by a
kernel boundary, and a change of cocycle by a boundary permits a
change of lift by a differential. These are the long exact sequence's
connecting maps and they commute with maps of coefficient
sequences. When the coefficient module is injective the groups in
positive degree vanish, by Proposition~\ref{res:A01}. Thus the
cochain long exact sequences have the universality proved by
dimension shifting in Theorem~\ref{res:R09}. The identity on
invariants extends uniquely to their comparison with the
right-derived long exact sequences. See
\cite[Sections 6.1 and 6.4]{weibel}.
\end{proof}

\section{Restriction, Shapiro, and group extensions}

\begin{theorem}[Shapiro's lemma]\label{res:A04}
If \(H\leq G\) and \(L\) is a left \(k[H]\)-module, put
\[
 \operatorname{Coind}_H^G L=\operatorname{Hom}_{k[H]}(k[G],L),
 \qquad (g\varphi)(x)=\varphi(xg).
\]
For every \(n\) there is a natural isomorphism
\(H^n(G,\operatorname{Coind}_H^GL)\simeq H^n(H,L)\). On
homogeneous group cochains the comparison is restriction to tuples
in \(H\), followed by evaluation at \(1\) on the coefficient.
\end{theorem}

\begin{proof}
For any \(G\)-module \(X\), a \(G\)-map from \(X\) to the coinduced
module has evaluation at \(1\) an \(H\)-map \(X|_H\to L\). Conversely
an \(H\)-map \(u\) gives \(x\mapsto(a\mapsto u(ax))\), first for
\(a\in G\) and then by \(k\)-linearity. This map has the required
\(H\)-linearity in \(a\) and \(G\)-linearity in \(x\) by the given
action. The two formulas are inverse and commute with both module
variables. The restriction of \(B_\bullet(G)\) to \(H\) is still a
free resolution since \(k[G]\) is free as a left \(k[H]\)-module
on representatives of \(H\backslash G\). Hence the adjunction applied
in every degree gives the computation on that resolution.
The inclusion \(B_\bullet(H)\to B_\bullet(G)|_H\) lifts the
identity of \(k\), so comparison of resolutions
\cite[Chapter 10]{p103} identifies it with the canonical comparison
on cohomology. Following evaluation at \(1\) by this comparison
restricts homogeneous tuples as claimed. The same maps of complexes
for all \(L\) prove naturality. See \cite[Theorem 6.3.2]{weibel}.
\end{proof}

\begin{theorem}[The Lyndon--Hochschild--Serre spectral sequence]
\label{res:A05}
Let \(N\) be normal in \(G\), \(Q=G/N\), and let \(M\) be a left
\(k[G]\)-module. There is a natural first-quadrant spectral sequence
\[
 E_2^{p,q}=H^p(Q,H^q(N,M))\ \Longrightarrow\ H^{p+q}(G,M)
\]
with finite target filtration in each degree. The \(Q\)-action on
the coefficient is induced from conjugation on \(N\) and the action on
\(M\). Its two edges are inflation and restriction, the latter
taking values in the invariants. In particular it gives
\begin{align*}
 0\longrightarrow H^1(Q,M^N)
  &\xrightarrow{\mathrm{inf}}H^1(G,M)
  \xrightarrow{\mathrm{res}} H^1(N,M)^Q\\
  &\xrightarrow{d_2}H^2(Q,M^N)
  \xrightarrow{\mathrm{inf}}H^2(G,M).
\end{align*}
\end{theorem}

\begin{proof}
The functor \(M\mapsto M^N\) takes values in \(Q\)-modules. An
element \(g\) sends an \(N\)-invariant element to an \(N\)-invariant
element by normality, and \(N\) acts trivially there. A \(G\)-map
from an inflated \(Q\)-module factors uniquely through \(M^N\),
which proves that this functor is right adjoint to inflation.
Inflation is exact on underlying modules. Its right adjoint
sends injectives to injectives: lifting along a monomorphism
before applying the adjunction is exactly lifting along its
inflation. In addition, restriction from \(G\) to \(N\) preserves
injectives because the induction functor \(k[G]\otimes_{k[N]}-\)
is an exact left adjoint, by freeness as a right \(k[N]\)-module.
Thus if \(M\to I^\bullet\) is a \(G\)-injective resolution, the
cohomology of \(C^\bullet=(I^\bullet)^N\) is \(H^q(N,M)\) as a
\(k\)-module, with a canonical \(Q\)-action.

For the identification of the action, on homogeneous \(N\)-cochains
use
\[
 (g\varphi)(n_0,\ldots,n_j)
 =g\,\varphi(g^{-1}n_0g,\ldots,g^{-1}n_jg).
\]
This respects \(N\)-equivariance and commutes with the deletion
differential. On degree-zero invariants it acts through the action
on the coefficient. For \(a\in N\), equivariance identifies this
cochain operation with pullback by simultaneous right translation
by \(a\) on homogeneous tuples. A homotopy on the resolution from
the identity to this translation is
\[
 (n_0,\ldots,n_j)\longmapsto
 \sum_{i=0}^{j}(-1)^i
 (n_0,\ldots,n_i,n_i a,\ldots,n_j a).
\]
Deletions before and after the repeated vertex cancel against the
corresponding prisms of the boundary faces. The two deletions at
each transition cancel between consecutive terms, except for the
two end faces, giving right translation minus identity. This homotopy
is \(N\)-linear. The action on cohomology thus factors through \(Q\).
It agrees with the injective-resolution action: use the
first-quadrant array of homogeneous cochains \(C^j(N,I^r)\).
Taking its cohomology first in the resolution direction produces
\(C^*(N,M)\), since the bar terms are free. Taking it first in the
bar direction produces \((I^\bullet)^N\), since the restricted
terms are injective. Each of the two natural augmentations is a
quasi-isomorphism by Theorem~\ref{res:B04}, and both are \(G\)-linear.

Invariants satisfy \((M^N)^Q=M^G\). The hypotheses of
Theorem~\ref{res:G04} therefore hold for \(M\mapsto M^N\) followed
by \(Q\)-invariants, and its object case gives the asserted page
and the finite abutment filtration. To identify the row edge,
apply derived \(Q\)-invariants to the map
\((M^N)[0]\to (I^\bullet)^N\). On bar cochains it is compared
with the inclusion of coefficients and pullback of tuples along
\(G\to Q\). In detail, for each \(I^r\) pullback gives
\[
 C^j(Q,(I^r)^N)\longrightarrow C^j(G,I^r).
\]
It commutes with both differentials. Both horizontal computations
collapse onto degree zero by injectivity of \(I^r\) and \((I^r)^N\).
There the map is the identity of \((I^r)^G\). The augmentations
from the \(M\)-cochains are compatible with these maps, which
identifies the row edge with the usual inflation map. The column
edge on the models is induced by the inclusion
\((I^\bullet)^G\to (I^\bullet)^N\), with cohomology in the fixed
part. This is the usual restriction: restriction of homogeneous
tuples gives a map \(C^j(G,I^r)\to C^j(N,I^r)\), which restricts
to precisely that inclusion on degree-zero invariants; the two
augmentations of the arrays again commute. The five terms then
follow from Corollary~\ref{res:G05}. See \cite[Section 6.8]{weibel}.
\end{proof}

\section{The Chevalley--Eilenberg resolution}

Let \(\mathfrak g\) be a Lie algebra over \(k\); its bracket is
alternating and satisfies Jacobi. Write \(U=U(\mathfrak g)\).
It is the quotient of the tensor algebra by the relations
\(xy-yx=[x,y]\). The augmentation kills \(\mathfrak g\); a left
\(U\)-module determines its underlying Lie module by restriction.
In this section the underlying \(k\)-module of \(\mathfrak g\) is
free, with no restriction on the rank. Write
\(\widehat{\boldsymbol x}_i\) or \(\widehat{\boldsymbol x}_{ij}\)
for the wedge formed by omitting the indicated entries, preserving
the order of those that remain.

\begin{proposition}[The Chevalley--Eilenberg complex]\label{res:A06}
On \(V_n=U\otimes_k \bigwedge^n\mathfrak g\) put
\begin{align*}
 \partial(u\otimes x_1\wedge\cdots\wedge x_n)
 &=\sum_{i=1}^n (-1)^{i-1}u x_i\otimes
                       \widehat{\boldsymbol x}_i\\
 &\quad+\sum_{1\leq i<j\leq n}(-1)^{i+j}u\otimes
                   [x_i,x_j]\wedge \widehat{\boldsymbol x}_{ij}.
\end{align*}
These maps form an augmented \(U\)-linear chain complex \(V_\bullet\to k\).
Under evaluation on \(1\otimes-\), the complex
\(\operatorname{Hom}_U(V_\bullet,M)\) has degree \(n\) term
\(C^n(\mathfrak g,M)=
 \operatorname{Hom}_k(\bigwedge^n\mathfrak g,M)\).
The coboundary is
\begin{align*}
 (df)(x_1,\ldots,x_{n+1})
 &=\sum_{i=1}^{n+1}(-1)^{i-1}x_i f(\widehat{\boldsymbol x}_i)\\
 &\quad+\sum_{1\leq i<j\leq n+1}(-1)^{i+j}
                 f([x_i,x_j],\widehat{\boldsymbol x}_{ij}).
\end{align*}
\end{proposition}

\begin{proof}
The two formulas for the boundary are multilinear and alternating:
the sign of a removal is the exterior sign of moving the removed
factor to the front, with an additional minus sign for removal of
a pair. If two entries are equal, terms omitting neither retain
them and vanish, and terms omitting just one cancel in pairs.
The bracket of the two equal entries is zero. Thus the formula
descends to exterior powers without a restriction on the
characteristic. Right multiplication of the \(U\)-factor in a
summand commutes with its left \(U\)-action, which gives
\(U\)-linearity. The augmentation kills the image of \(\partial:V_1\to V_0\).

For the square of the differential, collect terms by the original
indices involved. Move them to the front preserving the order of
every other factor. This multiplies both contributions of a given
type by the same exterior sign. For two removed entries \(x,y\),
the two multiplication operations give \(u(xy-yx)\); multiplication
of the newly created bracket gives \(-u[x,y]\). These cancel by
the enveloping relation. For a removed \(x\) and a disjoint
bracket \([y,z]\), the products obtained in the two orders are
\(-ux\otimes[y,z]\) and \(+ux\otimes[y,z]\), with the remaining
factors understood. Two disjoint brackets yield opposite terms
because reversing their orders reverses the order of the two
bracket factors in the wedge. For a triple \(x,y,z\) used in a
nested bracket, the contribution is
\[
 u\otimes\bigl([[x,y],z]-[[x,z],y]+[[y,z],x]\bigr),
\]
wedged with the unused factors. This vanishes by Jacobi. These
cases exhaust the two applications of the differential, proving
\(\partial^2=0\). Finally a \(U\)-linear map \(V_n\to M\) is
determined by its value on \(1\otimes-\), with the inverse given
by \(u\otimes w\mapsto u f(w)\). Precomposition by \(\partial\)
now gives the displayed coboundary, naturally in \(M\).
See \cite[Section 7.7]{weibel}.
\end{proof}

\paragraph{The ordered monomials in the enveloping algebra.}
Choose a totally ordered \(k\)-basis of \(\mathfrak g\). In a word
in basis elements, replace a descending adjacent pair \(yx\) by
\(xy+[y,x]\), expressing the bracket in the basis. This procedure
terminates: in each resulting monomial either the word length
decreases or the length is unchanged and the inversion count
decreases. A word contains only finitely many letters and each
bracket is a finite linear combination, irrespective of the size
of the basis.

Disjoint replacements commute. The only overlap of two different
applicable replacements is a strictly descending triple \(zyx\).
Starting with either pair and commuting the remaining descending
pairs gives the same length-three term \(xyz\). The difference
of the two expressions of smaller length is
\[
 [y,x]\,z-z\,[y,x]+y\,[z,x]-[z,x]\,y+[z,y]\,x-x\,[z,y].
\]
Each difference of concatenations of two basis-linear factors
reduces to their Lie bracket. The resulting sum is
\([[y,x],z]+[y,[z,x]]+[[z,y],x]=0\) by Jacobi. These checks apply
with any unused prefix and suffix. Inducting on length and inversion
count now gives unique linear normal forms. Indeed, to compare the
first reductions, commute them in the disjoint case and use the
computed overlap difference in the overlapping case, and normalize
the resulting terms using the induction hypothesis. The normal form
map kills every instance of the defining relation multiplied by
a prefix and suffix, and thus kills the ideal they generate.
Conversely every replacement is the defining relation in the quotient.
Ordered words therefore form a \(k\)-basis of \(U\), and the
associated graded algebra for the word-length filtration is
\(\operatorname{Sym}_k(\mathfrak g)\): the generators commute in
the top degree, and the sorted monomials on each side are the
same basis. This is the free-module case of the
Poincar\'e--Birkhoff--Witt theorem of \cite[Theorem 7.3.7]{weibel}.

\begin{theorem}[Chevalley--Eilenberg computation of Ext]\label{res:A07}
If \(\mathfrak g\) is free as a \(k\)-module, then the augmentation
of Proposition~\ref{res:A06} is a free \(U(\mathfrak g)\)-resolution
of \(k\). In particular the \(n\)-th cohomology of its cochain
complex is naturally \(\operatorname{Ext}^n_{U(\mathfrak g)}(k,M)\).
\end{theorem}

\begin{proof}
Filter \(V_n\) increasingly by total weight: weight at most \(a\)
means that \(u\) belongs to the word-length step \(F_{a-n}U\),
interpreted as zero for \(a-n<0\). The multiplication part of the
boundary preserves weight and the bracket part lowers it by one.
The sorted-word description above identifies the associated graded
complex with the Koszul chain complex on the variables of
\(\operatorname{Sym}_k(\mathfrak g)\); the differential is the sum
of signed deletions with multiplication by the deleted variable.
The augmentation on its degree-zero terms is the quotient map
by the variables.

For a finite subset of the basis, take the Koszul complex on that
subset over the whole polynomial ring. Starting with no variables,
adjoining one variable adds the two-term factor whose boundary
is multiplication by that variable. Equivalently the larger complex
is the chain mapping cone of multiplication by the new variable
on the smaller complex, with the tensor differential sign of
Chapter~\ref{ch:double-complexes}. The long exact sequence of that
cone computes its homology: by induction the smaller complex has
only degree-zero homology, the quotient by its variables, and
the new variable acts injectively on this polynomial quotient.
The cone thus has degree-zero homology the quotient by all variables
in the enlarged subset and zero homology in positive degrees.
For an infinite basis every wedge chain uses only a finite subset.
A positive-degree cycle is therefore a boundary in some such
finite-subset complex. The degree-zero homology of the union is
the quotient by all variables, namely \(k\). This argument proves
exactness of the augmented associated graded complex, in every
weight since the Koszul boundary preserves that weight.

Now let \(z\in V_n\) be a cycle for \(n>0\), or let \(z\in V_0\)
lie in the kernel of the augmentation. It has finite filtration
weight. Its leading class is a cycle of the augmented graded
complex. By graded exactness lift a graded antecedent to a
chain of one higher degree in the same filtration step. Subtracting
its boundary from \(z\) lowers its filtration weight. Repeating
terminates because the filtration in degree \(n\) is zero below
weight \(n\). The sum of the chosen antecedents shows that \(z\)
is a boundary. The augmentation is onto, and the modules \(V_n\)
are free over \(U\), since ordered wedges of basis elements form
a \(k\)-basis of \(\bigwedge^n\mathfrak g\).
The resolution computes Ext by Theorem~\ref{res:R11}, with the
identification of cochains in Proposition~\ref{res:A06}.
See \cite[Theorem 7.7.2 and Corollary 7.7.3]{weibel}.
\end{proof}

\begin{corollary}[Low-degree Lie algebra cohomology]\label{res:A08}
For the coefficients of Proposition~\ref{res:A06},
\[
 H^0(\mathfrak g,M)=M^{\mathfrak g},\qquad
 H^1(\mathfrak g,M)=
 \frac{\{f:\mathfrak g\to M\text{ \(k\)-linear}\mid
 f([x,y])=x f(y)-y f(x)\}}
 {\{x\mapsto x m\mid m\in M\}} .
\]
The corresponding differentials are \(dm(x)=xm\) and
\(df(x,y)=xf(y)-yf(x)-f([x,y])\).
\end{corollary}

\begin{proof}
These are the formulas of Proposition~\ref{res:A06} with zero
and one input in the cochain, respectively. For a degree-zero
vector \(m\), applying the second formula to \(x\mapsto xm\)
gives \(x(ym)-y(xm)-[x,y]m=0\), using the defining module
relation. Thus quotienting cocycles by these boundaries gives
exactly derivations modulo inner derivations. All identifications
are natural in \(M\), and Theorem~\ref{res:A07} identifies them
with the corresponding Ext groups. See
\cite[Definition 7.4.3 and Theorem 7.4.7]{weibel}.
\end{proof}

\section{The Lie algebra Hochschild--Serre comparison}

\begin{theorem}[An ideal and its quotient]\label{res:A09}
Let \(k\) be a field, \(\mathfrak g\) a Lie algebra over \(k\),
\(\mathfrak h\) an ideal, and \(\mathfrak q=\mathfrak g/\mathfrak h\).
For any \(\mathfrak g\)-module \(M\) there is a natural first-quadrant
spectral sequence
\[
 E_2^{p,q}=H^p(\mathfrak q,H^q(\mathfrak h,M))
           \ \Longrightarrow\ H^{p+q}(\mathfrak g,M)
\]
with finite target filtration in each degree. Its row and column
edges are induced by inflation and restriction of Chevalley--Eilenberg
cochains, and the restriction edge lands in the \(\mathfrak q\)-invariants.
The five terms in degree one and two are
\begin{align*}
 0\longrightarrow H^1(\mathfrak q,M^{\mathfrak h})
 &\xrightarrow{\mathrm{inf}}H^1(\mathfrak g,M)
 \xrightarrow{\mathrm{res}}H^1(\mathfrak h,M)^{\mathfrak q}\\
 &\xrightarrow{d_2}H^2(\mathfrak q,M^{\mathfrak h})
 \xrightarrow{\mathrm{inf}}H^2(\mathfrak g,M).
\end{align*}
\end{theorem}

\begin{proof}
Choose a basis of \(\mathfrak h\) and extend it to a basis of
\(\mathfrak g\). Order the complementary letters before the
\(\mathfrak h\)-letters. The ordered-word computation preceding
Theorem~\ref{res:A07} shows that the enveloping algebra of
\(\mathfrak h\) embeds in \(U(\mathfrak g)\), and that
\(U(\mathfrak g)\) is free as a right \(U(\mathfrak h)\)-module
on the ordered complementary words. Hence restriction preserves
injectives, being right adjoint to exact induction by tensoring
with that free module. The \(\mathfrak h\)-invariant part of a
\(\mathfrak g\)-module is a \(\mathfrak q\)-module:
if \(x\in\mathfrak h\), \(y\in\mathfrak g\), and \(xm=0\) for all
such \(x\), then \(x(ym)=y(xm)+[x,y]m=0\).
The invariant functor is right adjoint to exact inflation
because a map from an inflated module has image annihilated
by \(\mathfrak h\). It thus sends injectives to injectives.
The composite with \(\mathfrak q\)-invariants is precisely
\(\mathfrak g\)-invariants. All three categories of modules
have enough injectives. Theorems~\ref{res:A07} and~\ref{res:R09}
identify the derived invariants with the cohomology groups in
the statement, and Theorem~\ref{res:G04} applies.

For clarity, the quotient action may also be computed on the
cochains of Proposition~\ref{res:A06}. For \(y\in\mathfrak g\)
define
\[
 (L_y f)(x_1,\ldots,x_j)
 =y f(x_1,\ldots,x_j)-
   \sum_{i=1}^{j} f(x_1,\ldots,[y,x_i],\ldots,x_j).
\]
The bracket terms are in \(\mathfrak h\). The action property
follows from the module relation and Jacobi, which identifies the
commutator of the adjoint operators with the adjoint of the bracket.
It commutes with the cochain differential: on action terms this
uses \(y(xm)-x(ym)=[y,x]m\), and on bracket terms it uses that
\([y,-]\) is a derivation of the bracket. For \(a\in\mathfrak h\)
put \((\iota_a f)(x_1,\ldots,x_{j-1})=f(a,x_1,\ldots,x_{j-1})\),
and set \(\iota_a=0\) in degree zero. In
\(d\iota_a f+\iota_a df\) the action and bracket terms
on the \(x_i\) alone cancel in pairs after moving \(a\) to the
first argument. The terms left are the action by \(a\) and
one term \(-f(x_1,\ldots,[a,x_i],\ldots,x_j)\) per index.
Thus \(d\iota_a+\iota_a d=L_a\), so the action descends to
\(\mathfrak q\) on cohomology.

To compare it with the action used by derived invariants, choose
a \(\mathfrak g\)-injective resolution \(M\to I^\bullet\).
The double complex \(C^j(\mathfrak h,I^r)\) has the displayed
\(\mathfrak g\)-action by cochain maps. The same two augmentations
as in the group proof apply: exactness of
\(\operatorname{Hom}_k(\bigwedge^j\mathfrak h,-)\) computes the
resolution-direction cohomology, while restriction preserving
injectives makes the Chevalley--Eilenberg direction vanish in
positive degrees for each \(I^r\). By Theorem~\ref{res:B04},
the total cohomology comparisons are quasi-isomorphisms and
identify the two actions. This proves the needed identification on
the second page.

Finally, the quotient map \(\mathfrak g\to\mathfrak q\) pulls
back alternating cochains with values in \(M^{\mathfrak h}\),
and its pullback commutes with the differential because it is
a Lie map compatible with the action on these coefficients.
Restriction is likewise a cochain map using the ideal inclusion.
For each \(I^r\), pullback compares
\(C^*(\mathfrak q,(I^r)^{\mathfrak h})\) with
\(C^*(\mathfrak g,I^r)\). In degree-zero invariants it is the
identity on \((I^r)^{\mathfrak g}\); the higher cohomology on
both sides vanishes by injectivity. Thus it identifies the row
edge induced by \((M^{\mathfrak h})[0]\to (I^\bullet)^{\mathfrak h}\)
with inflation, because the comparison commutes with the coefficient
augmentations induced by \(M\to I^\bullet\). In the same way,
restriction of cochains
computes the inclusion of invariant complexes
\((I^\bullet)^{\mathfrak g}\to(I^\bullet)^{\mathfrak h}\)
giving the column edge. Apply Corollary~\ref{res:G05} to get
the five terms. See \cite[Section 7.5]{weibel}.
\end{proof}

\Needspace{24\baselineskip}
\section{Coinvariants and coefficient group homology}

\begin{theorem}[Group homology and its coefficient sequence]
\label{res:A10}
Let \(G\) be a discrete group, put \(R=k[G]\), and let \(M\) be
a left \(R\)-module. Its coinvariants are
\[
 M_G=M/\langle gm-m:g\in G,\ m\in M\rangle_k
       \simeq k\otimes_R M,
\]
where \(k\) is the trivial right \(R\)-module. For \(n\ge0\) define
\[
 H_n(G,M)=\operatorname{Tor}^{R}_n(k,M).
\]
These are naturally the left derived functors of coinvariants,
with \(H_0(G,M)=M_G\). They can be computed from the complex
\[
 C_n(G,M)=\bigoplus_{(g_1,\ldots,g_n)\in G^n}M .
\]
Write \(m[g_1|\cdots|g_n]\) for a vector in the displayed summand
and \(C_0(G,M)=M\). Its differential for \(n\ge1\) is
\begin{align*}
 \partial\bigl(m[g_1|\cdots|g_n]\bigr)
 &=g_1^{-1}m[g_2|\cdots|g_n]\\
 &\quad+\sum_{i=1}^{n-1}(-1)^i
       m[g_1|\cdots|g_i g_{i+1}|\cdots|g_n]\\
 &\quad+(-1)^n m[g_1|\cdots|g_{n-1}].
\end{align*}
Every short exact sequence \(0\to M'\to M\to M''\to0\) of
left \(R\)-modules gives a natural long exact sequence
\begin{align*}
 \cdots &\longrightarrow H_n(G,M')\longrightarrow H_n(G,M)
       \longrightarrow H_n(G,M'')\\
 &\xrightarrow{\delta}H_{n-1}(G,M')\longrightarrow\cdots
       \longrightarrow M_G\longrightarrow M''_G\longrightarrow0 .
\end{align*}
\end{theorem}

\begin{proof}
In \(k\otimes_R M\), the tensor relation identifies
\(1\otimes gm\) with \(1g\otimes m=1\otimes m\). Thus
\(m\mapsto1\otimes m\) factors through \(M_G\). Its inverse sends
\(a\otimes m\) to \(a[m]\) for \(a\in k\), and the augmentation
formula for the trivial right action makes this inverse balanced.
This proves the zeroth identification and shows that coinvariants
are right exact.

Take the free left bar resolution \(B_\bullet(G)\to k\) of
\ref{res:A02}. On its underlying \(k\)-modules put
\(b\cdot g=g^{-1}b\), extending \(k\)-linearly to a right
\(R\)-action. Inversion reverses products, so
\((b\cdot g)\cdot h=b\cdot(gh)\). The left-equivariant
differential is also right-equivariant for this action.
Call this augmented complex \(B_\bullet^{\mathrm r}(G)\to k\).
The inverse twist of a free left \(R\)-module is a free right
\(R\)-module: if its left basis is \((e_\alpha)\), then
the elements \(e_\alpha\cdot g=g^{-1}e_\alpha\) are the same
underlying \(k\)-basis as \(g e_\alpha\).
Its exactness is the exactness already proved for the underlying
bar complex, so this is a free right resolution of the trivial
right module.

Use the same free basis
\([g_1|\cdots|g_n]=(1,g_1,g_1g_2,\ldots,g_1\cdots g_n)\)
as in \ref{res:A02}, now as a right-module basis.
Then \(B_n^{\mathrm r}(G)\otimes_R M\) identifies with the
displayed direct sum by
\([g_1|\cdots|g_n]\otimes m\mapsto m[g_1|\cdots|g_n]\).
In the deletion boundary, the first face of this tuple is
the left translate \(g_1[g_2|\cdots|g_n]\).
In the twisted module it is
\([g_2|\cdots|g_n]\cdot g_1^{-1}\); balancing therefore puts
\(g_1^{-1}\) on the coefficient \(m\). The interior and final
faces are unchanged. This proves the differential formula,
and \(\partial^2=0\) also follows from the deletion calculation.
The exponent \(-1\) in its first term records precisely this
choice of right-module basis. At degree one the image is
spanned by \(g^{-1}m-m\), the same submodule as that spanned by
\(gm-m\), in agreement with the formula for \(H_0\).

To compare with the derived functor in the coefficient variable,
choose a projective left resolution \(P_\bullet\to M\).
The first-quadrant double complex
\(B_p^{\mathrm r}(G)\otimes_R P_q\) has finite diagonals.
Each \(B_p^{\mathrm r}(G)\) is free, so resolving \(M\) in the
\(q\)-direction identifies the total homology with the homology
of \(B_\bullet^{\mathrm r}(G)\otimes_R M\).
Each \(P_q\) is flat, so resolving \(k\) in the \(p\)-direction
identifies that same homology with the homology of
\(k\otimes_R P_\bullet\). These are the two resolution
comparisons of \ref{res:B05}. The latter complex computes the
left derived coinvariant functor by \ref{res:R07}; equivalently
place \(P_q\) in cochain degree \(-q\) as in \ref{res:R14}.
Thus the bar model, the two computations of Tor, and the derived
coinvariants agree naturally. Comparison maps between projective
resolutions are unique up to homotopy, so the comparisons are
independent of the choices.

Finally, the complex \(C_n(G,-)\) is a direct sum of copies of
the coefficient module, hence is exact in every degree.
A coefficient short exact sequence therefore gives a short
exact sequence of chain complexes. For a cycle \(z''\) in the
quotient complex, lift it to \(z\) in the middle complex.
The vector \(\partial z\) belongs to the image of the first
complex. Its unique preimage is a cycle by injectivity and
\(\partial^2=0\), and its homology class defines
\(\delta[z'']\). Changing the lift changes that preimage by a
boundary; changing the cycle representative by a boundary has
the same effect after choosing a lifted boundary. These maps
commute with maps of coefficient sequences and give the
displayed long exact sequence. To compare the connecting maps
with the coefficient-derived construction, take a horseshoe of
projective coefficient resolutions for the three modules.
The two augmentations of each double complex used above give
maps of short exact sequences of complexes: the horseshoe is
degreewise split, and the right bar terms are free. Both maps
induce the displayed homology isomorphisms. Naturality of the
lift-and-boundary construction transports the same connecting
map across both comparisons.
See \cite[Section 6.1 and Application 6.5.4]{weibel}.
\end{proof}
